\section{KV260 Accelerator Methodology}

\subsection{Verification-First Design Flow}

The hardware flow follows four levels. First, a host C++ reference reproduces
the exported integer logits for the supplied INT16 input. Second, each real
requantization ratio is approximated by an integer multiplier and shift and is
validated against the C++ reference on the declared validation set. Third, the
HLS top function and self-checking testbench are verified by C/RTL
co-simulation. Finally, the packaged design is integrated with AXI DMA on
KV260 and measured on board. C/RTL co-simulation validates generated RTL
against the C reference but does not replace board measurements
\cite{amd_hls_cosim}.

\subsection{Accelerator Variants}

The first design, H0, is a scalar functional baseline. H1 pipelines the
temporal convolution loops and stores static weights in on-chip memories. H2
introduces a controlled unroll factor for pointwise convolutions. H3 optionally
uses layer-level dataflow with explicit intermediate buffers. All variants use
the same checkpoint, quantization constants, test vector, clock target, and
measurement protocol. Therefore, changes in latency or resources can be
attributed to dataflow and parallelization rather than a changed model.

\subsection{Measurement Boundaries}

Two timing boundaries are reported: kernel-only and host-plus-DMA. Batch size
is one window because the intended streaming use case emits a new two-second
window every second. At least 1,000 inferences are used to report median and
95th-percentile latency. Power methodology, sampling interval, and idle-power
subtraction are documented. Post-route LUT, FF, BRAM, URAM, DSP, and achieved
clock frequency are taken from the implementation reports.
